
% ================================================================
% NEW MATERIAL FROM STANDALONE (ordered chiral homology paper)
% E_3 identification, explicit operations, chiral P_3 bracket,
% filtered E_3-chiral algebra, CFG comparison
% Inserted losslessly; macros to be adapted for memoir class
% ================================================================

  All non-canonically isomorphic to
  $(k+h^\vee) H^3(\fg)[[k+h^\vee]]$.
\end{enumerate}
\end{theorem}

\begin{theorem}[Deformation families of $\Ethree$-algebras]
\label{thm:e3-identification}
Let $\fg$ be a simple finite-dimensional Lie algebra and
$V_k(\fg)$ the affine Kac--Moody vertex algebra at
level~$k$.
Writing $\lambda = k + h^\vee$ for the departure from
critical level, the derived chiral centre
$Z^{\mathrm{der}}_{\mathrm{ch}}(V_k(\fg))$ of
Proposition~\textup{\ref{prop:e3-structure}} and the
CFG $\Ethree$-algebra $\cA^\lambda$ of
Theorem~\textup{\ref{thm:cfg}} are isomorphic as formal
deformation families of $\Ethree$-algebras:
\begin{equation}\label{eq:e3-families}
  \{Z^{\mathrm{der}}_{\mathrm{ch}}(V_k(\fg))
  \}_{k + h^\vee \in \lambda\CC[[\lambda]]}
  \;\simeq\; \{\cA^\lambda\}_{\lambda
  \in \lambda H^3(\fg)[[\lambda]]}
\end{equation}
over the base $\lambda H^3(\fg)[[\lambda]]$
\textup{(}with $\lambda = k + h^\vee$\textup{)}.
\begin{enumerate}[label=\textup{(\roman*)}]
\item \textup{(Matching deformation spaces.)}
  Both families are parametrised by
  $(k+h^\vee) H^3(\fg)[[k+h^\vee]]$:
  for the derived chiral centre, the deformation
  parameter is the departure from critical level;
  for the CFG algebra, it is the Chern--Simons coupling.
\item \textup{(Base point mismatch.)}
  At the critical level $k = -h^\vee$, the derived
  chiral centre $Z^{\mathrm{der}}_{\mathrm{ch}}
  (V_{-h^\vee}(\fg))$ is related to the
  Feigin--Frenkel centre
  $\mathrm{Fun}(\mathrm{Op}_{\fg^\vee}(D))$
  \textup{(}opers on the formal disk\textup{)}, not
  directly to $C^*(\fg)$. The isomorphism is
  perturbative: it holds in the formal neighbourhood
  of critical level, order by order in
  $\lambda = k + h^\vee$.
\item \textup{(Order-by-order uniqueness.)}
  For simple $\fg$, $H^3(\fg) \cong \CC$ is
  one-dimensional, so the deformation base is
  one-dimensional at each order in $\lambda$. The
  isomorphism~\eqref{eq:e3-families} is forced by
  the matching of associated graded algebras
  \textup{(}both $C^*(\fg)$\textup{)},
  matching $\Pthree$ brackets on the formal disk
  \textup{(}Theorem~\textup{\ref{thm:chiral-e3-cfg}(ii)}%
  \textup{)}, and the one-dimensionality of the
  obstruction space at each order.
\item \textup{(Sugawara constraint.)}
  The topological enhancement
  \textup{(}Theorem~\ref{thm:e3-cs}(iv)\textup{)} holds
  at generic $k$ but fails at $k = -h^\vee$, where
  the Sugawara element is undefined. The perturbative
  identification holds
  order by order in $(k + h^\vee)$, since the Sugawara
  element exists at any nonzero departure from critical
  level.
\end{enumerate}
\end{theorem}

\begin{proof}
The argument is obstruction-theoretic: we construct an
isomorphism of $\Ethree$-algebras order by order in
$\lambda = k + h^\vee$, using the one-dimensionality of
$H^3(\fg)$ for simple~$\fg$ to force uniqueness at each
stage.

\smallskip
\noindent\textit{Step 1: setup.}
Both $Z^{\mathrm{der}}_{\mathrm{ch}}(V_k(\fg))$ and
$\cA^\lambda$ are filtered $\Ethree$-algebras deforming
$C^*(\fg) = \Sym(\fg^\vee[-1])$ over the base
$\lambda H^3(\fg)[[\lambda]]$, with $\lambda = k + h^\vee$
(Theorem~\ref{thm:e3-cs}(ii) and
Theorem~\ref{thm:cfg}(iv) respectively).
Both filtrations are complete and exhaustive with
associated graded $C^*(\fg)$.
The deformation theory of $\Ethree$-algebras with fixed
associated graded $C^*(\fg)$ is controlled by the
$\Ethree$-deformation complex, whose relevant cohomology
group is $H^3(\fg)$: the space of $\Ethree$-deformations
of $C^*(\fg)$ at each order $p \geq 1$ in $\lambda$,
modulo those at order $< p$, is classified by
$H^3(\fg)$ (Lurie~\cite{HA},
Costello--Francis--Gwilliam~\cite{CFG26},
Theorem~1.4).

\smallskip
\noindent\textit{Step 2: induction on order.}
We construct the isomorphism
$\varphi \colon Z^{\mathrm{der}}_{\mathrm{ch}}(V_k(\fg))
\xrightarrow{\sim} \cA^\lambda$ as a compatible family of
isomorphisms $\varphi_p$ modulo $\lambda^{p+1}$, by
induction on~$p$.

\emph{Base case ($p = 0$):} both associated graded
algebras are $C^*(\fg)$
(Theorem~\ref{thm:chiral-e3-cfg}(i)), so
$\varphi_0 = \mathrm{id}_{C^*(\fg)}$.

\emph{Inductive step:} suppose $\varphi_p$ is an
$\Ethree$-isomorphism modulo $\lambda^{p+1}$.
The obstruction to extending $\varphi_p$ to an
isomorphism modulo $\lambda^{p+2}$ lies in the
space of $\Ethree$-deformations of $C^*(\fg)$ at
order $p+1$, which is $H^3(\fg)$. For simple~$\fg$,
$H^3(\fg) = \CC$ is one-dimensional, generated by the
Cartan three-form
$\omega_3(X, Y, Z) = (X, [Y, Z])$.
Two $\Ethree$-deformations of $C^*(\fg)$
that agree modulo $\lambda^{p+1}$ therefore differ at order
$p+1$ by a scalar multiple of $\omega_3$.

It remains to show that both constructions select the
\emph{same} scalar at each order. On the formal disk
$D = \Spec \CC[[z]]$, the restriction functor
$\Gamma(D, -)$ is symmetric monoidal and preserves
$\Ethree$-structures
(Theorem~\ref{thm:chiral-e3-cfg}, Step~1 of the
proof). By Theorem~\ref{thm:chiral-e3-cfg}(ii), the
chiral $\Pthree$ bracket restricted to~$D$ recovers
the CFG $\Pthree$ bracket at every order. Since the
$\Pthree$ bracket determines the $\Ethree$-deformation
class (the bracket on generators
$\fg^*[-1] \otimes \fg^*[-1] \to \CE^{\mathrm{ch}}(\fg_k)$
encodes the full deformation datum for biderivations, by
the argument of
Lemma~\ref{lem:bv-p3-commutativity}, Step~2), and the
space of $\fg$-equivariant corrections at each order is
one-dimensional (spanned by the Killing-form term
$k\,(a,b)\,\partial$), the two constructions agree at
order $p+1$. The isomorphism $\varphi_{p+1}$ extends
$\varphi_p$.

\smallskip
\noindent\textit{Step 3: passage to the limit.}
The compatible family $\{\varphi_p\}_{p \geq 0}$
determines an isomorphism
$\varphi = \varprojlim_p\,\varphi_p$
of filtered $\Ethree$-algebras over
$\lambda\,H^3(\fg)[[\lambda]]$, completing the proof
of~\eqref{eq:e3-families}.
Part~(iv) is immediate: the topological enhancement
(Theorem~\ref{thm:e3-cs}(iv)) requires the Sugawara
element, which exists for $\lambda \neq 0$ and is
compatible with the isomorphism~$\varphi$ order by order
(the Sugawara element at order~$p$ depends only on the
data modulo~$\lambda^{p+1}$, where the two families
already agree).
\end{proof}

% ----------------------------------------------------------------
\subsection{Explicit $\Ethree$ operations for
$V_k(\mathfrak{sl}_2)$}
\label{subsec:e3-explicit}

\subsection{Explicit $\Ethree$ operations for
$V_k(\mathfrak{sl}_2)$}
\label{subsec:e3-explicit}

The abstract $\Ethree$ structure of
Proposition~\ref{prop:e3-structure} acquires concrete content
for $\fg = \mathfrak{sl}_2$: the five-dimensional derived
chiral centre $\CC \oplus \mathfrak{sl}_2[-1] \oplus \CC[-2]$
carries an explicit cup product, Gerstenhaber bracket,
$\Pthree$ bracket, and BV operator, all computable from the
OPE.

\medskip
\noindent
\textbf{The underlying graded vector space.}
By Proposition~\ref{prop:e3-structure} and the chiral
Hochschild computation of
Example~\ref{ex:km}, the derived chiral centre of
$V_k(\mathfrak{sl}_2)$ at generic level
$k \neq -2$ is the graded vector space
\begin{equation}\label{eq:e3-graded-space}
  Z^{\mathrm{der}}_{\mathrm{ch}}(V_k(\mathfrak{sl}_2))
  \;=\;
  \underbrace{\CC}_{=\, \HH^0}
  \;\oplus\;
  \underbrace{\mathfrak{sl}_2[-1]}_{=\, \HH^1}
  \;\oplus\;
  \underbrace{\CC[-2]}_{=\, \HH^2},
\end{equation}
with total dimension $1 + 3 + 1 = 5$. Write $\mathbf{1}$
for the generator of $\HH^0 = \CC$ (the vacuum),
$\{e, f, h\}$ for the standard basis of
$\HH^1 = \mathfrak{sl}_2$
(outer derivations: infinitesimal level-preserving
deformations parametrised by $\fg$ itself), and $\eta$
for the generator of $\HH^2 = \CC$
(the central charge deformation cocycle). The
$\mathfrak{sl}_2$-module structure is: $\HH^0$ trivial,
$\HH^1$ adjoint, $\HH^2$ trivial.

\medskip
\noindent
\textbf{Conventions.}
The Killing form normalisation is
$(e, f) = (f, e) = 1$, $(h, h) = 2$, all other pairings zero.
Structure constants: $[e, f] = h$, $[h, e] = 2e$,
$[h, f] = -2f$. The inverse Casimir is
$\Omega = e \otimes f + f \otimes e
+ \tfrac{1}{2}\,h \otimes h$, and $h^\vee = 2$.
The KZ coupling constant is
$h_{\mathrm{KZ}} = 1/(k + h^\vee) = 1/(k + 2)$.

\begin{remark}[Convention for $\hbar$]
\label{rem:hbar-convention}
Throughout this paper, $\hbar$ denotes the
\emph{KZ coupling constant}
\begin{equation}\label{eq:hbar-convention}
  \hbar \;=\; h_{\mathrm{KZ}}
  \;=\; \frac{1}{k + h^\vee},
\end{equation}
the normalisation in the KZ connection
$\nabla = d + \hbar \sum \Omega_{ij}\,dz_{ij}$.
It vanishes as $k \to \infty$ (weak-coupling limit) and
diverges at the critical level $k = -h^\vee$.
In the $\mathfrak{sl}_2$ computations, $h^\vee = 2$,
so $\hbar = 1/(k+2)$.
The \emph{departure from critical level} is the inverse
quantity $k + h^\vee = 1/\hbar$; we always write this
explicitly as $k + h^\vee$ rather than introducing a
separate symbol. Critical level corresponds to
$k + h^\vee = 0$ (equivalently, $\hbar \to \infty$);
weak coupling corresponds to $k + h^\vee \to \infty$
(equivalently, $\hbar \to 0$).
For the quantum group parameter: $q = e^{2\pi i \hbar}$.
\end{remark}

An $\Ethree$-algebra carries three levels of
operations:
\begin{itemize}
\item Level $0$ ($\Einf$): a graded-commutative product
  $\mu\colon \HH^p \otimes \HH^q \to \HH^{p+q}$.
\item Level $1$ ($\Etwo$ beyond $\Eone$): the
  Gerstenhaber bracket
  $[-,-]\colon \HH^p \otimes \HH^q \to \HH^{p+q-1}$
  of degree~$-1$, derived from the fundamental class of
  $S^1 = \Conf_2(\RR^2)/\RR^+$.
\item Level $2$ ($\Ethree$ beyond $\Etwo$): the $\Pthree$
  bracket
  $\{-,-\}\colon \HH^p \otimes \HH^q \to \HH^{p+q-2}$
  of degree~$-2$, derived from the fundamental class of
  $S^2 = \Conf_2(\RR^3)/\RR^+$.
\end{itemize}
We compute each in turn.

\begin{proposition}[Explicit $\Ethree$ operations on
$Z^{\mathrm{der}}_{\mathrm{ch}}(V_k(\mathfrak{sl}_2))$]
\label{prop:e3-explicit-sl2}
Let $k \neq -2$.
The $\Ethree$-algebra structure on
$Z^{\mathrm{der}}_{\mathrm{ch}}(V_k(\mathfrak{sl}_2))$
is determined by the following data.

\smallskip
\noindent
\textup{(I) Cup product} ($\Einf$ level).
\begin{enumerate}[label=\textup{(\alph*)}]
\item $\mathbf{1}$ is the unit:
  $\mathbf{1} \cdot f = f \cdot \mathbf{1} = f$ for all $f$.
\item
  $\mu(X, Y) = 0$ for all $X, Y \in \HH^1 = \mathfrak{sl}_2$.
\item
  $\mu(X, \eta) = 0$ for all $X \in \HH^1$, and
  $\eta \cdot \eta = 0$.
\end{enumerate}

\smallskip
\noindent
\textup{(II) Gerstenhaber bracket} ($\Etwo$ level,
degree~$-1$).
\begin{gather*}
  [\mathbf{1}, -] = 0, \qquad
  [X, Y] = 0 \quad \text{for all }
  X, Y \in \HH^1 = \mathfrak{sl}_2, \\
  [X, \eta] = 0 \quad \text{for all } X \in \HH^1,
  \qquad [\eta, \eta] = 0.
\end{gather*}

\smallskip
\noindent
\textup{(III) The $\Pthree$ bracket} ($\Ethree$ level,
degree~$-2$).
\begin{align}
  \label{eq:e3-p3-killing}
  \{X, Y\}
  &= h_{\mathrm{KZ}} \cdot (X, Y)
  = \frac{1}{k + h^\vee}\,(X, Y)
  \quad \text{for } X, Y \in \HH^1,
  \\
  \label{eq:e3-p3-derivation}
  \{X, \eta\}
  &= 0
  \quad \text{for } X \in \HH^1,
  \\
  \label{eq:e3-p3-unit}
  \{\mathbf{1}, -\}
  &= 0,
  \\
  \label{eq:e3-p3-top}
  \{\eta, \eta\}
  &= 0,
\end{align}
where $(X, Y)$ denotes the normalised Killing form
and $h_{\mathrm{KZ}} = 1/(k + h^\vee) = 1/(k + 2)$ is
the KZ coupling constant
\textup{(}Remark~\textup{\ref{rem:hbar-convention}}\textup{)}.
\end{proposition}

\begin{proof}
We establish each level by $\mathfrak{sl}_2$-equivariance,
degree counting, and the BV relation.

\textit{Cup product.}
The cup product $\mu\colon \HH^p \otimes \HH^q
\to \HH^{p+q}$ is graded-commutative:
$\mu(a, b) = (-1)^{pq}\,\mu(b, a)$.
For $X, Y \in \HH^1$, this gives
$\mu(X, Y) = -\mu(Y, X)$: the restriction to
$\HH^1 \otimes \HH^1 \to \HH^2 = \CC$ is an
$\mathfrak{sl}_2$-invariant antisymmetric bilinear form
on the adjoint representation valued in the trivial
representation.
Since
$\Lambda^2(\mathfrak{sl}_2)^{\mathfrak{sl}_2}
\cong \Hom_{\mathfrak{sl}_2}(\Lambda^2(\mathrm{ad}), \CC)
= 0$
(the exterior square of the adjoint decomposes as
$\Lambda^2(\mathrm{ad}) \cong \mathrm{ad}$,
which has no trivial summand), the cup product
$\mu(X, Y) = 0$ for all $X, Y \in \HH^1$.

The products $\mu(\HH^1, \HH^2) \subset \HH^3 = 0$ and
$\mu(\HH^2, \HH^2) \subset \HH^4 = 0$ vanish for degree
reasons.

\textit{BV operator.}
The framed $\Etwo$ structure on $\BarSig(V_k(\mathfrak{sl}_2))$
(from genus-$1$ sewing) equips the derived centre with a BV
operator $\Delta\colon \HH^n \to \HH^{n-1}$ satisfying
$\Delta^2 = 0$ and the BV relation
\begin{equation}\label{eq:bv-rel-sl2}
  [f, g]
  = \Delta(f \cdot g) - (\Delta f) \cdot g
  - (-1)^{|f|}\,f \cdot (\Delta g).
\end{equation}

The component $\Delta\colon \HH^1 \to \HH^0$ is an
$\mathfrak{sl}_2$-equivariant linear map from the adjoint
representation to the trivial representation. Since the
adjoint is irreducible and non-trivial,
$\Delta|_{\HH^1} = 0$.

The component $\Delta\colon \HH^2 \to \HH^1$ is an
$\mathfrak{sl}_2$-equivariant linear map from the trivial
to the adjoint. Again, $\Delta|_{\HH^2} = 0$.

All other components vanish trivially
($\Delta|_{\HH^0}$ lands in $\HH^{-1} = 0$;
$\Delta|_{\HH^n} = 0$ for $n \geq 3$). Therefore
$\Delta = 0$ on the entire derived centre.

\textit{Gerstenhaber bracket.}
With $\Delta = 0$ and $\mu|_{\HH^1 \otimes \HH^1} = 0$,
the BV relation~\eqref{eq:bv-rel-sl2} gives:
for $X, Y \in \HH^1$,
\[
  [X, Y]
  = \Delta(\mu(X, Y)) - (\Delta X) \cdot Y
  + X \cdot (\Delta Y)
  = \Delta(0) - 0 + 0 = 0.
\]
For $X \in \HH^1$, $\eta \in \HH^2$:
$[X, \eta] = \Delta(\mu(X, \eta)) - (\Delta X) \cdot \eta
- (-1)^1\,X \cdot (\Delta \eta)
= 0 - 0 + 0 = 0$
(since $\mu(X, \eta) \in \HH^3 = 0$,
$\Delta X = 0$, $\Delta\eta = 0$).
Similarly $[\eta, \eta] \in \HH^3 = 0$.

The vanishing of the Gerstenhaber bracket on the derived
centre is a structural consequence of the
$\mathfrak{sl}_2$-equivariance constraints: the adjoint
representation provides no equivariant maps to or from the
trivial representation. This is specific to the
\emph{derived centre} $\HH^*(\BarSig, \BarSig)$; the
Gerstenhaber bracket on the full Hochschild cochain complex
$C^*(\BarSig, \BarSig)$ is non-trivial (it carries the
commutator of derivations at the cochain level).

\textit{$\Pthree$ bracket.}
The degree-$(-2)$ bracket
$\{-,-\}\colon \HH^p \otimes \HH^q \to \HH^{p+q-2}$
satisfies the graded symmetry
\begin{equation}\label{eq:p3-symmetry}
  \{a, b\}
  = -(-1)^{(|a|-2)(|b|-2)}\,\{b, a\}
\end{equation}
of a degree-$(-2)$ Lie bracket, together with the Leibniz rule
$\{a, b \cdot c\}
= \{a, b\} \cdot c + (-1)^{(|a|-2)|b|}\,b \cdot \{a, c\}$.

\textit{Step 1: $\{\mathbf{1}, -\} = 0$.}
The Leibniz rule applied to
$\mathbf{1} = \mathbf{1} \cdot \mathbf{1}$ gives
$\{\mathbf{1}, f\} = 2\,\{\mathbf{1}, f\}$
for all $f$; hence $\{\mathbf{1}, f\} = 0$.

\textit{Step 2: $\{X, Y\} = h_{\mathrm{KZ}}\,(X, Y)$.}
For $X, Y \in \HH^1$,
$\{X, Y\} \in \HH^0 = \CC$.
The symmetry~\eqref{eq:p3-symmetry} at
$|X| = |Y| = 1$ gives
$\{X, Y\} = -(-1)^{(-1)(-1)}\,\{Y, X\}
= -(-1)^1\,\{Y, X\}
= \{Y, X\}$:
the bracket is \emph{symmetric} on $\HH^1 \otimes \HH^1$.
By $\mathfrak{sl}_2$-equivariance, $\{X, Y\}$ is an
invariant symmetric bilinear form on the adjoint
representation valued in $\CC$. The space of such forms is
one-dimensional, spanned by the Killing form $(X, Y)$.
Hence $\{X, Y\} = \alpha \cdot (X, Y)$ for a scalar
$\alpha$ depending on $k$.

The scalar $\alpha$ is determined by the $r$-matrix
residue, independently of any comparison with the CFG
construction. The $\Pthree$ bracket on the derived
chiral centre $Z^{\mathrm{der}}_{\mathrm{ch}}(V_k(\fg))
= \HH^*(\BarSig(V_k(\fg)), \BarSig(V_k(\fg)))$ is the
$\Etwo$-Hochschild cohomology operation derived from
the fundamental class of $S^2 \subset \Conf_2(\RR^3)$.
At degree~$2$, this operation is the residue of the
classical $r$-matrix paired with the fundamental class
of $S^1 \subset \Conf_2(\CC)$. In the KZ normalisation
(see~\eqref{eq:kz-sl2-degree2} and the conventions
of~\S\ref{subsec:sl2-chiral-e3}):
% AP126/AP148: KZ convention r(z) = Omega/((k+h^v)z).
% At k=0, r = Omega/(h^v z) != 0 (non-abelian; correct for KZ).
% At k=-h^v, r diverges (Sugawara singularity).
\begin{equation}\label{eq:p3-from-r-matrix-residue}
  r(z) = \frac{\Omega}{(k + h^\vee)\,z},
  \qquad
  \operatorname{Res}_{z=0}\bigl[r(z)\bigr]
  = \frac{\Omega}{k + h^\vee}.
\end{equation}
The residue acts on $X \otimes Y \in \fg \otimes \fg$
by contraction with the inverse Casimir:
$\Omega(X \otimes Y) = (X, Y)$ (the normalised Killing
form). The $\Pthree$ bracket on $\HH^1 \cong \fg$ is
therefore
\begin{equation}\label{eq:p3-bracket-normalisation}
  \{X, Y\}
  = \operatorname{Res}_{z=0}
  \bigl[r(z)(X \otimes Y)\bigr]
  = \frac{(X, Y)}{k + h^\vee}
  = h_{\mathrm{KZ}} \cdot (X, Y),
\end{equation}
giving $\alpha = h_{\mathrm{KZ}} = 1/(k + h^\vee)$.
For $\mathfrak{sl}_2$, $h^\vee = 2$ and
$\alpha = 1/(k + 2)$.

Verification: at $k \to \infty$,
$h_{\mathrm{KZ}} \to 0$ and the bracket vanishes (the
weak-coupling limit is the abelian $\Ethree$). At the
critical level $k = -h^\vee$, $h_{\mathrm{KZ}}$ diverges
(the bracket degenerates, reflecting the Feigin--Frenkel
singularity). Both limits are consistent with
the behaviour of the KZ connection.

\textit{Step 3: $\{X, \eta\} = 0$.}
The bracket $\{X, \eta\} \in \HH^1 = \mathfrak{sl}_2$
for $X \in \HH^1$, $\eta \in \HH^2$ is an
$\mathfrak{sl}_2$-equivariant map
$\mathrm{ad} \otimes \CC \to \mathrm{ad}$,
hence $\{X, \eta\} = \beta\,X$ for some scalar $\beta$.
The $\Pthree$ Jacobi identity applied to
$X, Y \in \HH^1$ and $\eta \in \HH^2$,
combined with
$\{X, Y\} = h_{\mathrm{KZ}}\,(X, Y) \in \HH^0$, gives
$\{\{X, \eta\}, \eta\} = 0$ (since
$\{\{X, \eta\}, \eta\} \in \HH^{-1} = 0$) and
$\{\eta, \{X, \eta\}\} \in \HH^{-1} = 0$,
while $\{X, \{\eta, \eta\}\} = \{X, 0\} = 0$.
The Jacobi identity is thus satisfied for all $\beta$;
the constraint comes instead from the self-consistency
of the $\Ethree$ deformation.
At the cohomological level, one applies the iterated Jacobi
identity to $X, Y \in \HH^1$ and $\eta$:
\begin{align*}
  \{\{X, Y\}, \eta\}
  &= \{X, \{Y, \eta\}\}
  - (-1)^{(|X|-2)(|Y|-2)}\,\{Y, \{X, \eta\}\} \\
  &= \{X, \beta Y\} - (-1)^1\,\{Y, \beta X\} \\
  &= \beta\,h_{\mathrm{KZ}}\,(X, Y)
  + \beta\,h_{\mathrm{KZ}}\,(Y, X) \\
  &= 2\beta\,h_{\mathrm{KZ}}\,(X, Y).
\end{align*}
The left side:
$\{\{X, Y\}, \eta\}
= \{h_{\mathrm{KZ}}(X, Y)\,\mathbf{1}, \eta\}
= h_{\mathrm{KZ}}(X, Y)\,\{\mathbf{1}, \eta\} = 0$.
Hence $2\beta\,h_{\mathrm{KZ}}\,(X, Y) = 0$ for all
$X, Y$. Since
$(e, f) = 1 \neq 0$ and $h_{\mathrm{KZ}} \neq 0$ at
generic $k$, we conclude $\beta = 0$.

\textit{Step 4: $\{\eta, \eta\} = 0$.}
The bracket $\{\eta, \eta\} \in \HH^2 = \CC$ and
the symmetry~\eqref{eq:p3-symmetry} at
$|\eta| = 2$ gives
$\{\eta, \eta\} = -(-1)^{0 \cdot 0}\,\{\eta, \eta\}
= -\{\eta, \eta\}$, hence $\{\eta, \eta\} = 0$.
\end{proof}

\begin{remark}[The Killing form as the $\Ethree$ shadow
of the Cartan $3$-form]
\label{rem:cartan-3-form}
The bracket~\eqref{eq:e3-p3-killing} has a transparent
topological origin. The $\Ethree$-deformation space is
$(k+h^\vee)\,H^3(\mathfrak{sl}_2)[[k+h^\vee]] \cong (k+h^\vee)\,\CC[[k+h^\vee]]$,
one-dimensional at each order in $(k+h^\vee)$
(Theorem~\ref{thm:cfg}(iv)). The generator of
$H^3(\mathfrak{sl}_2)$ is the Cartan $3$-form
\begin{equation}\label{eq:cartan-3form}
  \omega_3(X, Y, Z)
  = (X, [Y, Z]),
\end{equation}
the unique (up to scalar) $\mathfrak{sl}_2$-invariant element
of $\Lambda^3(\mathfrak{sl}_2^*)$. The degree-$(-2)$ bracket
$\{X, Y\}$ is the $\Ethree$-operation induced by this
deformation class scaled by
$\hbar = 1/(k + h^\vee)$:
\[
  \{X, Y\}
  = \hbar \cdot (X, Y).
\]
The mechanism is \emph{not} literal contraction of the
antisymmetric $3$-form $\omega_3$: the $\Pthree$ bracket
is symmetric on $\HH^1$ (by the graded
symmetry~\eqref{eq:p3-symmetry} at $|X| = |Y| = 1$),
while any contraction of $\omega_3$ in two slots is
antisymmetric. The connection is indirect:
the one-dimensional deformation space
$H^3(\mathfrak{sl}_2)$, generated by $\omega_3$,
parametrises $\Ethree$-deformations of $C^*(\fg)$
(Theorem~\ref{thm:cfg}(iv));
on cohomology, $\mathfrak{sl}_2$-equivariance forces
the induced $\Pthree$ bracket on
$\HH^1 \cong \mathfrak{sl}_2$ to be proportional to
the Killing form $(X, Y)$, the unique invariant
symmetric bilinear form. The coefficient
$h_{\mathrm{KZ}} = 1/(k + h^\vee)$ is fixed by the
$r$-matrix residue
(equation~\eqref{eq:p3-from-r-matrix-residue}), not by
comparison with CFG.
The $\Pthree$ bracket on $Z^{\mathrm{der}}_{\mathrm{ch}}$
is the binary shadow of the ternary Cartan deformation class,
with the Killing form appearing through equivariance rather
than through direct contraction.
\end{remark}

% ----------------------------------------------------------------
\subsection{Quantum corrections from the vertex $R$-matrix:
$\Ethree$ operations for the EK quantum VOA}
\label{subsec:e3-ek-quantum}

The classical $V_k(\mathfrak{sl}_2)$ has trivial braiding
$S(z) = \id$; the Etingof--Kazhdan quantum VOA
$V_{\mathrm{EK}}$ of Example~\ref{ex:ek-qvoa} has
$S(z) = \id + h_{\mathrm{KZ}}\,\Omega/z + O(z^{-2})$. Does
the non-trivial $R$-matrix change the $\Ethree$-algebra
structure on the derived chiral centre? On cohomology, no:
the $\Pthree$ bracket coincides with the classical answer
(Proposition~\ref{prop:e3-explicit-sl2}). On cochains, yes:
the quantum corrections encode the braided monoidal structure
of $\mathrm{Rep}(Y_\hbar(\mathfrak{sl}_2))$.

\medskip
\noindent
\textbf{The underlying graded vector space.}
The EK quantum VOA is a flat formal deformation of
$V_k(\mathfrak{sl}_2)$ over $\CC[[h_{\mathrm{KZ}}]]$, with
$h_{\mathrm{KZ}} = 1/(k+2)$.
Since the chiral Hochschild cohomology is computed by
a bounded complex of finite-dimensional
$\mathfrak{sl}_2$-modules, the deformation does not
change the underlying graded vector space:
\begin{equation}\label{eq:ek-e3-graded}
  Z^{\mathrm{der}}_{\mathrm{ch}}(V_{\mathrm{EK}})
  \;=\;
  \CC[[h_{\mathrm{KZ}}]]
  \;\oplus\;
  \mathfrak{sl}_2[[h_{\mathrm{KZ}}]][-1]
  \;\oplus\;
  \CC[[h_{\mathrm{KZ}}]][-2],
\end{equation}
with total rank $5$ over $\CC[[h_{\mathrm{KZ}}]]$.
The generators
$\mathbf{1} \in \HH^0$, $\{e, f, h\} \in \HH^1$, and
$\eta \in \HH^2$ are defined as in
Proposition~\ref{prop:e3-explicit-sl2}, now taken
over the base ring $\CC[[h_{\mathrm{KZ}}]]$.

\medskip
\noindent
\textbf{Equivariance under the quantum $R$-matrix.}
The Yang $R$-matrix $R(u) = u\,I + \hbar\,P$ on
$V \otimes V$ is $\mathfrak{sl}_2$-equivariant:
$[R(u),\, \rho(X) \otimes I + I \otimes \rho(X)] = 0$
for all $X \in \mathfrak{sl}_2$.
The full vertex $R$-matrix
$S(z)$ on $V_{\mathrm{EK}}^{\otimes 2}$
(equation~\eqref{eq:ek-vertex-rmatrix}) is likewise
$\mathfrak{sl}_2$-equivariant, since it is constructed from
$\mathcal{R}(z)$, the spectral universal $R$-matrix of
the Yangian $Y_\hbar(\mathfrak{sl}_2)$, which
commutes with the $\mathfrak{sl}_2$-action by the
defining property of the Drinfeld coproduct.

This equivariance is the key structural constraint: all
three levels of the $\Ethree$ operations on
$Z^{\mathrm{der}}_{\mathrm{ch}}(V_{\mathrm{EK}})$ are
$\mathfrak{sl}_2$-equivariant maps between representations.
The equivariance arguments of
Proposition~\ref{prop:e3-explicit-sl2} therefore apply
verbatim to the quantum case.

\begin{proposition}[{$\Ethree$ operations on
$Z^{\mathrm{der}}_{\mathrm{ch}}(V_{\mathrm{EK}})$}]
\label{prop:e3-ek-quantum}
Let $V_{\mathrm{EK}}$ be the Etingof--Kazhdan quantum
vertex algebra for $\mathfrak{sl}_2$ at KZ coupling
$h_{\mathrm{KZ}} = 1/(k+2) \neq 0$
\textup{(}Example~\textup{\ref{ex:ek-qvoa}}\textup{)}.
The $\Ethree$-algebra structure on the derived chiral
centre $Z^{\mathrm{der}}_{\mathrm{ch}}(V_{\mathrm{EK}})$
coincides with the classical $\Ethree$ structure of
Proposition~\textup{\ref{prop:e3-explicit-sl2}}: the
cup product and Gerstenhaber bracket vanish, and the
$\Pthree$ bracket is
\begin{equation}\label{eq:ek-p3-bracket}
  \{X, Y\}_q \;=\; h_{\mathrm{KZ}} \cdot (X, Y)
  \qquad
  \textup{for } X, Y \in \HH^1 = \mathfrak{sl}_2,
\end{equation}
with all other brackets zero. In particular, the
vertex $R$-matrix $S(z)$ does not renormalise the
$\Pthree$ bracket on cohomology.
\end{proposition}

\begin{proof}
The argument has four ingredients: equivariance,
one-dimensionality of the obstruction space,
exactness of the rational classical $r$-matrix, and a
$\GRT_1$ consistency check.

\textit{Step 1: cup product and BV operator.}
The cup product $\mu \colon \HH^p \otimes \HH^q
\to \HH^{p+q}$ is a graded-commutative product on
$\mathfrak{sl}_2$-modules over $\CC[[h_{\mathrm{KZ}}]]$.
For $X, Y \in \HH^1
= \mathfrak{sl}_2[[h_{\mathrm{KZ}}]]$,
the product $\mu(X, Y) \in \HH^2
= \CC[[h_{\mathrm{KZ}}]]$
is an $\mathfrak{sl}_2$-equivariant antisymmetric
bilinear form on the adjoint representation:
$\Lambda^2(\mathrm{ad})^{\mathfrak{sl}_2} = 0$,
so $\mu(X, Y) = 0$ over $\CC[[h_{\mathrm{KZ}}]]$.

The BV operator $\Delta \colon \HH^n \to \HH^{n-1}$
is $\mathfrak{sl}_2$-equivariant. The components
$\Delta|_{\HH^1} \colon \mathrm{ad} \to \CC$ and
$\Delta|_{\HH^2} \colon \CC \to \mathrm{ad}$ vanish
by Schur's lemma. Hence $\Delta = 0$ on the derived
centre, and the BV relation gives
$[X, Y] = 0$ for all $X, Y \in \HH^1$, just as in the
classical case.

\textit{Step 2: the $\Pthree$ bracket is proportional
to the Killing form.}
For $X, Y \in \HH^1$, the bracket
$\{X, Y\}_q \in \HH^0 = \CC[[h_{\mathrm{KZ}}]]$ is an
$\mathfrak{sl}_2$-equivariant symmetric bilinear
form on the adjoint representation:
$\Sym^2(\mathrm{ad})^{\mathfrak{sl}_2} \cong \CC$.

(the notation $\kappa$ here denotes the Koszul invariant,
not the Killing form).
At critical level $k = -h^\vee$ (i.e.\ $k + h^\vee = 0$), the
curvature vanishes and
$\CE^{\mathrm{ch}}(\fg_{-h^\vee}) = (\Sym^c(\fg^*[-1]) \otimes
\omega_X,\, d_{\CE})$
is an uncurved commutative ($\Einf$) factorisation
coalgebra on~$X$.
\end{definition}

\begin{construction}[The chiral $\Pthree$ bracket]
\label{constr:chiral-p3-bracket}
The PVA $\lambda$-bracket on $V_k(\fg)$ is
(see \cite{Lorgat26I}, \S5.11):
\begin{equation}\label{eq:pva-bracket-recall}
  \{J^a \mathbin{{}_\lambda} J^b\}
  \;=\;
  f^{ab}{}_c\, J^c \;+\; k\,(a,b)\,\lambda,
\end{equation}
where $f^{ab}{}_c$ are the structure constants of $\fg$ and
$(a,b)$ is the Killing form.

At the level of the chiral Chevalley--Eilenberg complex, the
$\lambda$-bracket induces a degree-$(-2)$ bracket
\begin{equation}\label{eq:chiral-p3-def}
  \{-,-\}^{\mathrm{ch}}
  \colon
  \CE^{\mathrm{ch}}(\fg_k)
  \otimes
  \CE^{\mathrm{ch}}(\fg_k)
  \;\longrightarrow\;
  \CE^{\mathrm{ch}}(\fg_k),
\end{equation}
defined on generators $\phi_a = (J^a)^* \in \fg^*[-1]$ by
the two components of the PVA bracket:
\begin{enumerate}[label=\textup{(\roman*)}]
\item \textup{(Lie bracket term.)}
  The structure constant term $f^{ab}{}_c\,J^c$ contributes
  the Chevalley--Eilenberg differential $d_{\CE}$ at the
  algebra level. On the dual (coalgebra) side, it induces a
  coderivation.
  This is the classical $\Pthree$ bracket of the
  uncurved complex $\CE^{\mathrm{ch}}(\fg_0)$: the degree-$(-2)$
  Lie bracket
  \begin{equation}\label{eq:p3-lie-component}
    \{\phi_a, \phi_b\}^{\mathrm{ch}}_{\mathrm{Lie}}
    \;=\;
    f^{ab}{}_c\,\phi_c.
  \end{equation}
  This bracket is the coinvariant shadow of the ordered bar
  $R$-matrix at the critical level.
\item \textup{(Cocycle term.)}
  The level-dependent term $k\,(a,b)\,\lambda$ contributes a
  degree-$(-2)$ symmetric pairing valued in $\cO_X$:
  \begin{equation}\label{eq:p3-cocycle-component}
    \{\phi_a, \phi_b\}^{\mathrm{ch}}_{\mathrm{cocycle}}
    \;=\;
    k\,(a,b) \cdot \partial,
  \end{equation}
  where $\partial$ denotes the $\cD$-module action
  (the spectral parameter~$\lambda$ of the PVA bracket
  becomes the derivation $\partial_z$ in the $\cD$-module
  structure). This term is the cocycle that controls
  the deformation away from the commutative
  ($\Einf$) factorisation structure.
\end{enumerate}

The total chiral $\Pthree$ bracket is
\begin{equation}\label{eq:chiral-p3-total}
  \{\phi_a, \phi_b\}^{\mathrm{ch}}
  \;=\;
  f^{ab}{}_c\,\phi_c
  \;+\;
  k\,(a,b)\,\partial.
\end{equation}
\end{construction}

\begin{proposition}[The chiral $\Pthree$ structure]
\label{prop:chiral-p3-structure}
The bracket~\eqref{eq:chiral-p3-total} equips the chiral
Chevalley--Eilenberg complex $\CE^{\mathrm{ch}}(\fg_k)$
with the structure of a $\Pthree$-algebra object in
$\cD\textrm{-}\mathrm{mod}(X)$. Concretely:
\begin{enumerate}[label=\textup{(\roman*)}]
\item The bracket is a degree-$(-2)$ Lie bracket satisfying
  graded antisymmetry and the Jacobi identity.
  The Jacobi identity follows from the Jacobi identity of the
  PVA $\lambda$-bracket, which in turn encodes the Jacobi
  identity of the Lie algebra $\fg$ together with the
  $\fg$-invariance of the Killing form.
\item The bracket is compatible with the commutative product
  on $\Sym^c(\fg^*[-1])$ via the Leibniz rule:
  \[
    \{a, b \cdot c\}^{\mathrm{ch}}
    = \{a, b\}^{\mathrm{ch}} \cdot c
    + (-1)^{(|a|-2)|b|}\, b \cdot \{a, c\}^{\mathrm{ch}}.
  \]
\item The bracket is $\cD$-linear in the sense that it
  respects the flat connection on
  $\CE^{\mathrm{ch}}(\fg_k)$ inherited from the
  $\cD$-module structure: the chiral $\Pthree$ bracket
  is a morphism of $\cD$-modules, not merely of
  $\cO_X$-modules.
\end{enumerate}
\end{proposition}

\begin{proof}
The PVA $\lambda$-bracket on $V_k(\fg)$ satisfies the PVA
axioms (sesquilinearity, Leibniz, skew-symmetry, Jacobi);
these are the $\lambda$-bracket translations of the classical
master equation for the Poisson structure.
The identification of the $\lambda$-bracket with a
degree-$(-2)$ Lie bracket on the shifted dual follows the
general correspondence between PVA brackets and shifted
Poisson structures (Beilinson--Drinfeld~\cite{BD04},
\S3.8): a PVA $\lambda$-bracket of conformal weight $1$
(where $\lambda$ has weight $1$) corresponds to a
$\Pthree$ bracket (degree $-2$ = $-\dim_\CC X - 1$, where
$\dim_\CC X = 1$) on the $\cD$-module of chiral cochains.

The $\cD$-linearity follows because the PVA
$\lambda$-bracket is a $\cD$-module morphism by construction:
the derivation $\partial$ acts via $\lambda \mapsto \lambda$
in the $\lambda$-bracket formalism, and this action is
compatible with the Leibniz rule by the sesquilinearity
axiom of the PVA.
\end{proof}

% ----------------------------------------------------------------
\subsection{The filtered $\Ethree$-chiral algebra}
\label{subsec:filtered-e3-chiral}

The chiral $\Pthree$ bracket and the chiral
Chevalley--Eilenberg complex assemble into a single object:
a filtered $\Ethree$-chiral algebra
$\CE^{\mathrm{ch}}_k(\fg)$, a factorisation $\cD$-module
on $\Ran(X)$ whose global sections on the formal disk
recover the CFG $\Ethree$-algebra.

\begin{definition}[The filtered $\Ethree$-chiral algebra]
\label{def:filtered-e3-chiral}
The \emph{chiral $\Ethree$-algebra}
\begin{equation}\label{eq:chiral-e3-algebra}
  \CE^{\mathrm{ch}}_k(\fg)
  \;:=\;
  \bigl(\CE^{\mathrm{ch}}(\fg_k),\;
  \{-,-\}^{\mathrm{ch}},\;
  \cD\textrm{-module structure}\bigr)
\end{equation}
is the chiral Chevalley--Eilenberg complex
(Definition~\ref{def:chiral-ce-complex}) equipped with
the chiral $\Pthree$ bracket
(Construction~\ref{constr:chiral-p3-bracket}) and the
$\cD$-module structure from the factorisation algebra
on~$X$. It carries a $(k+h^\vee)$-adic filtration
$F^\bullet \CE^{\mathrm{ch}}_k(\fg)$ defined by:
\begin{equation}\label{eq:hbar-filtration}
  F^p \CE^{\mathrm{ch}}_k(\fg)
  \;=\;
  (k+h^\vee)^p \cdot \CE^{\mathrm{ch}}_k(\fg).
\end{equation}
\end{definition}

\begin{theorem}[Structure of the chiral $\Ethree$-algebra]
\label{thm:chiral-e3-structure}
Let $\fg$ be a simple Lie algebra.
The chiral $\Ethree$-algebra
$\CE^{\mathrm{ch}}_k(\fg)$ satisfies:
\begin{enumerate}[label=\textup{(\roman*)}]
\item \textup{(Associated graded.)}
  The associated graded with respect to the $(k+h^\vee)$-adic
  filtration is the uncurved chiral Chevalley--Eilenberg
  complex at the critical level:
  \begin{equation}\label{eq:assoc-graded}
    \mathrm{gr}_F\,\CE^{\mathrm{ch}}_k(\fg)
    \;\simeq\;
    \CE^{\mathrm{ch}}(\fg_{-h^\vee})
    \;=\;
    \bigl(\Sym^c(\fg^*[-1]) \otimes \omega_X,\; d_{\CE}\bigr),
  \end{equation}
  which is an $\Einf$-chiral algebra \textup{(}commutative
  factorisation $\cD$-module on $\Ran(X)$\textup{)}.
\item \textup{(The $\Ethree$ structure.)}
  The chiral $\Pthree$ bracket
  \textup{(}Proposition~\textup{\ref{prop:chiral-p3-structure}}%
  \textup{)} deforms the $\Einf$ structure to
  a filtered $\Ethree$-chiral algebra:
  $\Etwo$ from the curve geometry
  \textup{(}Warning~\textup{\ref{warn:e1-vs-e2-source}}%
  \textup{)} $+ 1$ from the chiral $\Pthree$ bracket
  $= \Ethree$.
\item \textup{(Curvature and the deformation parameter.)}
  The curvature
  $m_0 = \tfrac{k + h^\vee}{2h^\vee}\,\kappa$ of the chiral
  Chevalley--Eilenberg complex is the obstruction to the
  $\Einf$ structure: at $k = -h^\vee$, the curvature vanishes
  and the algebra is commutative; at $k \neq -h^\vee$, the
  $\Pthree$ bracket and the curvature together deform the
  commutative structure to a filtered $\Ethree$-chiral
  algebra.
  The deformation parameter is $(k+h^\vee) \in
  (k+h^\vee)\,H^3(\fg)[[k+h^\vee]]$,
  matching the deformation space of
  Theorem~\textup{\ref{thm:e3-cs}(ii)}.
\item \textup{(Factorisation property.)}
  The $\Ethree$-chiral algebra
  $\CE^{\mathrm{ch}}_k(\fg)$ is a factorisation algebra on
  $X$ in the sense that the restriction to disjoint
  opens $U_1 \sqcup U_2 \hookrightarrow X$ induces
  \begin{equation}\label{eq:factorisation}
    \CE^{\mathrm{ch}}_k(\fg)\big|_{U_1 \sqcup U_2}
    \;\simeq\;
    \CE^{\mathrm{ch}}_k(\fg)\big|_{U_1}
    \otimes
    \CE^{\mathrm{ch}}_k(\fg)\big|_{U_2},
  \end{equation}
  where the tensor product is in the derived category of
  $\cD$-modules. The $\Ethree$ structure is compatible with
  the factorisation isomorphism.
\end{enumerate}
\end{theorem}
The following lemma supplies the key compatibility between
the BV operator $\Delta$ and the chiral $\Pthree$ bracket
used in the proof of part~(ii). The ``independent
geometric data'' heuristic (genus-$1$ sewing vs.\
genus-$0$ collision) motivates the result, but
the proof requires a precise argument.

\begin{lemma}[Commutativity of the BV operator and the
chiral $\Pthree$ bracket]
\label{lem:bv-p3-commutativity}
Let $\fg$ be a simple Lie algebra and $k \neq -h^\vee$.
The BV operator $\Delta$ on the derived chiral centre
$Z^{\mathrm{der}}_{\mathrm{ch}}(V_k(\fg))$
\textup{(}from the framed $\Etwo$ structure,
genus-$1$ sewing\textup{)} and the chiral $\Pthree$
bracket $\{-,-\}^{\mathrm{ch}}$
\textup{(}Construction~\textup{\ref{constr:chiral-p3-bracket}},
genus-$0$ collision residue\textup{)} satisfy
\begin{equation}\label{eq:bv-p3-commute}
  [\Delta, \{a,-\}^{\mathrm{ch}}] = 0
  \quad \text{for all } a \in \CE^{\mathrm{ch}}(\fg_k).
\end{equation}
\end{lemma}

\begin{proof}
The abstract $\Ethree$-algebra structure on the derived
chiral centre $Z^{\mathrm{der}}_{\mathrm{ch}}(V_k(\fg))$
is provided by Proposition~\ref{prop:e3-structure}(ii)
via the Higher Deligne Conjecture applied to the
$\Etwo$-coalgebra $\BarSig(V_k(\fg))$.
This abstract $\Ethree$ structure includes, in particular,
an abstract BV operator $\Delta_{\mathrm{HDC}}$
(from the framed $\Etwo$ structure) and an abstract
$\Pthree$ bracket $\{-,-\}_{\mathrm{HDC}}$
(the degree-$(-2)$ Lie bracket from the $S^2$ fundamental
class of $\Conf_2(\RR^3)/\RR^+$). These satisfy all
$\Ethree$ compatibility conditions, including
$[\Delta_{\mathrm{HDC}}, \{a, -\}_{\mathrm{HDC}}] = 0$,
by the operadic structure of $\Ethree$.

It remains to identify the explicit chiral $\Pthree$
bracket $\{-,-\}^{\mathrm{ch}}$ from
Construction~\ref{constr:chiral-p3-bracket} with
the abstract bracket $\{-,-\}_{\mathrm{HDC}}$. We argue
in three steps.

\textit{Step 1: matching on the associated graded.}
At $k = -h^\vee$ (the associated graded of the
$(k+h^\vee)$-adic filtration), the chiral
Chevalley--Eilenberg complex becomes the uncurved
commutative coalgebra $\CE^{\mathrm{ch}}(\fg_0)$,
which is an $\Einf$-chiral algebra.
Both $\{-,-\}^{\mathrm{ch}}$ and $\{-,-\}_{\mathrm{HDC}}$
reduce to the classical $\Pthree$ bracket on
$\Sym^c(\fg^*[-1]) \otimes \omega_X$ induced by the Lie
bracket of $\fg$: the bracket
$\{\phi_a, \phi_b\} = f^{ab}{}_c\,\phi_c$ on generators,
extended as a biderivation.
This classical bracket is uniquely determined by the
$\fg$-equivariance and degree constraints on
$\Sym^c(\fg^*[-1])$: the space of degree-$(-2)$
Lie brackets on $\Sym^c(\fg^*[-1])$ that are
biderivations of the commutative product and
$\fg$-equivariant is one-dimensional (generated by the
bracket induced by $\fg$ itself). Therefore
$\{-,-\}^{\mathrm{ch}}$ and $\{-,-\}_{\mathrm{HDC}}$
agree on $\mathrm{gr}_F$.

\textit{Step 2: matching order by order.}
The $(k+h^\vee)$-adic filtration on
$\CE^{\mathrm{ch}}(\fg_k)$ induces a filtration on the
space of degree-$(-2)$ Lie biderivations. At each order
$p \geq 1$ in $(k+h^\vee)$, the bracket is determined by
its value on generators
$\phi_a \otimes \phi_b \in \fg^*[-1] \otimes \fg^*[-1]$
(since both brackets are biderivations of $\mu$). On
generators, the space of $\fg$-equivariant $\cD$-linear
pairings $\fg^*[-1] \otimes \fg^*[-1] \to
\CE^{\mathrm{ch}}(\fg_k)$ of degree~$-2$ is spanned by
$f^{ab}{}_c\,\phi_c$ and $k\,(a,b)\,\partial$ (the
structure constants and the Killing form, the only
$\fg$-invariant tensors of the appropriate type for
simple~$\fg$). Step~1 fixes the structure-constant coefficient
$f^{ab}{}_c\,\phi_c$ at order~$0$.  At each subsequent
order $p \geq 1$, the \emph{correction} to the bracket on
generators lies in the subspace spanned by the Killing-form
term $k\,(a,b)\,\partial$ alone (since the
structure-constant component is already determined by
the lower orders and the biderivation property).
For simple~$\fg$, this correction space is
one-dimensional: $\Hom_\fg(\Sym^2\fg, \CC) \cong \CC$
(the Killing form is the unique $\fg$-invariant symmetric
bilinear form).  Both $\{-,-\}^{\mathrm{ch}}$ and
$\{-,-\}_{\mathrm{HDC}}$ are $\fg$-equivariant
$\cD$-linear biderivations with the same
order-$0$ value (Step~1), so their corrections at each
order $p$ coincide by the one-dimensionality of the
Killing-form component.  Therefore the two brackets agree
at each order.

\textit{Step 3: conclusion.}
Since $\{-,-\}^{\mathrm{ch}} = \{-,-\}_{\mathrm{HDC}}$
as formal power series in $(k+h^\vee)$, and
$[\Delta_{\mathrm{HDC}}, \{a,-\}_{\mathrm{HDC}}] = 0$
holds by the $\Ethree$ operadic structure,
the commutativity~\eqref{eq:bv-p3-commute} follows.
\end{proof}

\begin{remark}[Chain-level status of
Lemma~\ref{lem:bv-p3-commutativity}]
\label{rem:bv-p3-chain-level}
The proof reduces the commutativity to the identification
of the explicit chiral $\Pthree$ bracket with the abstract
HDC bracket. This identification is proved order by order
in $(k+h^\vee)$ using the one-dimensionality of the
deformation space, which holds for simple~$\fg$.
For non-simple~$\fg$, the deformation space
$H^3(\fg)[[k+h^\vee]]$ may be higher-dimensional, and the
order-by-order uniqueness argument requires a more refined
comparison.

The ``independent geometric data'' intuition (that
$\Delta$ comes from genus-$1$ sewing strata of
$\overline{\cM}_{1,n}$ while
$\{-,-\}^{\mathrm{ch}}$ comes from genus-$0$ collision
strata of $\overline{\cM}_{0,n}$) captures the operadic
reason: in the modular envelope of the $\Ethree$ operad,
genus-$1$ and genus-$0$ operations compose without
interference because they correspond to disjoint cells
in the modular operad. Making this precise at the
chain level in the chiral (as opposed to topological)
setting is what the identification argument above
accomplishes for simple~$\fg$.
\end{remark}

\begin{proof}[Proof of Theorem~\ref{thm:chiral-e3-structure}]
Part (i) is immediate from the filtration: the $(k+h^\vee)$-adic
filtration is exhaustive, and at $k = -h^\vee$ the curvature
and the cocycle term of the $\Pthree$ bracket both vanish,
leaving the uncurved commutative chiral coalgebra
$\CE^{\mathrm{ch}}(\fg_{-h^\vee})$.

Part (ii) requires three things: that the chiral
$\Pthree$ bracket is a degree-$(-2)$ Lie bracket
(Jacobi), that it satisfies the Leibniz rule with respect
to the commutative product (Poisson compatibility), and
that it is compatible with the $\Etwo$ structure from
the curve geometry, so that the combined data assembles
into an $\Ethree$-algebra. We prove each in turn.

\textit{Step 1: the $\Pthree$ Jacobi identity from the
PVA Jacobi.}
The PVA $\lambda$-bracket on $V_k(\fg)$ satisfies the
$(-1)$-shifted Jacobi identity
(see~\cite{Lorgat26II}, Appendix~A, Proposition~A.5):
\begin{equation}\label{eq:pva-jacobi-recall}
  \{a \mathbin{{}_\lambda} \{b \mathbin{{}_\mu} c\}\}
  - (-1)^{(|a|+1)(|b|+1)}\,
  \{b \mathbin{{}_\mu} \{a \mathbin{{}_\lambda} c\}\}
  = \{\{a \mathbin{{}_\lambda} b\}
  \mathbin{{}_{{\lambda+\mu}}} c\}.
\end{equation}
The chiral $\Pthree$ bracket
$\{-,-\}^{\mathrm{ch}}$
on $\CE^{\mathrm{ch}}(\fg_k)$ is obtained from the
$\lambda$-bracket by specialising $\lambda \mapsto \partial$
(the $\cD$-module derivation) and passing to the chiral
Chevalley--Eilenberg complex. On generators
$\phi_a, \phi_b, \phi_c \in \fg^*[-1]$
(cohomological degree $1$, shifted bracket-degree
$|\phi_a| - 2 = -1$), the degree-$(-2)$ Jacobi identity
is:
\begin{equation}\label{eq:p3-jacobi}
  \{\phi_a, \{\phi_b, \phi_c\}^{\mathrm{ch}}\}^{\mathrm{ch}}
  + \text{graded cyclic}
  = 0.
\end{equation}
To derive this from~\eqref{eq:pva-jacobi-recall}, set
$\lambda = \mu = 0$ (the constant-coefficient specialisation
that extracts the $\Pthree$ bracket from the
$\lambda$-bracket). The left side
of~\eqref{eq:pva-jacobi-recall} becomes
$\{\phi_a, \{\phi_b, \phi_c\}^{\mathrm{ch}}\}^{\mathrm{ch}}
- (-1)^{(|\phi_a|+1)(|\phi_b|+1)}\,
\{\phi_b, \{\phi_a, \phi_c\}^{\mathrm{ch}}\}^{\mathrm{ch}}$,
and the right side becomes
$\{\{\phi_a, \phi_b\}^{\mathrm{ch}},
\phi_c\}^{\mathrm{ch}}$ (at $\lambda + \mu = 0$).
The PVA $\lambda$-bracket has degree $-1$ (it is
$(-1)$-shifted); the chiral $\Pthree$ bracket has
degree~$-2$, accounting for the additional degree shift
from the $\cD$-module derivation $\partial$ (conformal
weight~$1$). The Koszul sign
$(-1)^{(|\phi_a|+1)(|\phi_b|+1)}$ in the PVA identity
translates to $(-1)^{(|\phi_a|-2)(|\phi_b|-2)}$ in the
$\Pthree$ identity after incorporating the degree shift
$|\,{-}\,| \mapsto |\,{-}\,| - 2$ appropriate to a
degree-$(-2)$ bracket. On $\fg^*[-1]$ where
$|\phi_a| = 1$, both signs evaluate to
$(-1)^{(-1)(-1)} = -1$. Rearranging
gives~\eqref{eq:p3-jacobi}.

For general elements of $\Sym^c(\fg^*[-1])$ (not just
generators), the Jacobi identity extends by the Leibniz
rule (Step~2 below): the bracket is a biderivation of the
free commutative algebra, so a trilinear identity that
holds on generators holds on all elements.

\textit{Step 2: the Leibniz rule from PVA sesquilinearity.}
The PVA right Leibniz rule
(see~\cite{Lorgat26II}, Appendix~A, Proposition~A.6)
states:
\[
  \{a \mathbin{{}_\lambda} (b \cdot c)\}
  = \{a \mathbin{{}_\lambda} b\} \cdot c
  + (-1)^{(|a|+1)|b|}\,
  b \cdot \{a \mathbin{{}_\lambda} c\}
  + \int_0^\lambda
  \{\{a \mathbin{{}_\mu} b\}
  \mathbin{{}_\lambda} c\}\,d\mu.
\]
At $\lambda = 0$, the integral term vanishes
(the integration range collapses to a point), leaving:
\[
  \{a, b \cdot c\}^{\mathrm{ch}}
  = \{a, b\}^{\mathrm{ch}} \cdot c
  + (-1)^{(|a|+1)|b|}\,
  b \cdot \{a, c\}^{\mathrm{ch}}.
\]
Adjusting the Koszul sign from the $(-1)$-shifted
convention to the degree-$(-2)$ convention:
$(|a|+1)|b|$ becomes $(|a|-2)|b|$ after the degree shift,
recovering the $\Pthree$ Leibniz rule
\[
  \{a, b \cdot c\}^{\mathrm{ch}}
  = \{a, b\}^{\mathrm{ch}} \cdot c
  + (-1)^{(|a|-2)|b|}\,
  b \cdot \{a, c\}^{\mathrm{ch}},
\]
as stated in
Proposition~\ref{prop:chiral-p3-structure}(ii). This
identifies the chiral $\Pthree$ bracket as a derivation of
the commutative product
$\mu$ on $\Sym^c(\fg^*[-1]) \otimes \omega_X$: a
commutative algebra equipped with a compatible
degree-$(-2)$ Lie bracket is a $\Pthree$-algebra.

\textit{Step 3: compatibility with $\Etwo$ and the
$\Ethree$ assembly.}
The $\Etwo$-coalgebra structure on $\BarSig$
(Warning~\ref{warn:e1-vs-e2-source}) and its framed
variant (from genus-$1$ sewing) equip the derived centre
with a BV operator $\Delta$ satisfying
$[f, g] = \Delta(f \cdot g) - (\Delta f) \cdot g
- (-1)^{|f|}\,f \cdot (\Delta g)$
(the BV relation~\eqref{eq:bv-rel-sl2}).
The $\Pthree$ bracket of degree~$-2$ supplements
this $\Etwo$ structure with one additional degree of
operadic freedom:
the $\Pthree$ bracket
(degree $-2 = -(3-1)$) corresponds to the fundamental
class of $S^2 = \Conf_2(\RR^3)/\RR^+$, the extra
generator needed for $\Ethree$ beyond~$\Etwo$.

The compatibility condition is that the $\Pthree$ bracket
must be a derivation of the Gerstenhaber bracket
$[-,-]$ of degree~$-1$ from the $\Etwo$ structure.
Since the Gerstenhaber bracket is determined by $\Delta$
and $\mu$ via the BV relation, it suffices to show that
$\{-,-\}^{\mathrm{ch}}$ commutes with $\Delta$ and is a
derivation of $\mu$ (already proved in Step~2). The
commutativity $[\Delta, \{a,-\}^{\mathrm{ch}}] = 0$
follows from Lemma~\ref{lem:bv-p3-commutativity}.
Given this, for any $a, b, c \in
\CE^{\mathrm{ch}}(\fg_k)$:
\begin{align*}
  \{a, [b, c]\}^{\mathrm{ch}}
  &= \{a, \Delta(b \cdot c)
  - (\Delta b) \cdot c
  - (-1)^{|b|}\,b \cdot (\Delta c)\}^{\mathrm{ch}} \\
  &= \Delta\{a, b \cdot c\}^{\mathrm{ch}}
  - \{a, (\Delta b) \cdot c\}^{\mathrm{ch}}
  - (-1)^{|b|}\,\{a, b \cdot (\Delta c)\}^{\mathrm{ch}},
\end{align*}
where the first term uses
$[\Delta, \{a,-\}^{\mathrm{ch}}] = 0$. Expanding each
term by the Leibniz rule of Step~2 and collecting via the
BV relation yields
\[
  \{a, [b, c]\}^{\mathrm{ch}}
  = [\{a, b\}^{\mathrm{ch}}, c]
  + (-1)^{(|a|-2)|b|}\,[b, \{a, c\}^{\mathrm{ch}}]:
\]
the $\Pthree$ bracket is a derivation of the
Gerstenhaber bracket.

This is the mixed Leibniz rule that characterises
$\Ethree = \Etwo + \Pthree$: the $\Pthree$ bracket acts
as a derivation of the $\Etwo$ Gerstenhaber bracket. By
the recognition principle for $\mathsf{E}_n$-algebras
(Lurie~\cite{HA}, Theorem~5.1.2.2;
Francis~\cite{Francis2013}, Theorem~4.1): a commutative
algebra with a compatible $\Etwo$ structure and a
degree-$(-2)$ Lie bracket satisfying the mutual Leibniz
rules (Jacobi, Poisson compatibility with the product, and
derivation of the Gerstenhaber bracket) is an
$\Ethree$-algebra. The three verifications above establish
all the required relations.

Part (iii) combines two observations. The curvature
$m_0 = \tfrac{k + h^\vee}{2h^\vee}\,\kappa$ is the bar-complex
encoding of the level departure (the identification
$\CE^{\mathrm{ch}}(\fg_k) \simeq B^{\mathrm{ch}}(V_k(\fg))$).
The cocycle term
$k\,(a,b)\,\partial$ of the $\Pthree$ bracket is linear
in~$k$ and hence linear in $(k + h^\vee)$.
Both vanish at $k = -h^\vee$, so the
deformation is controlled by $(k + h^\vee)$. The identification
of the deformation space as $(k+h^\vee) H^3(\fg)[[k+h^\vee]]$
then follows from the computation of
Theorem~\ref{thm:e3-cs}(ii): $H^3(\fg)$ is one-dimensional
for simple $\fg$, and the Cartan $3$-form
$\omega_3(X, Y, Z) = (X, [Y,Z])$ is the generator.

Part (iv) is the factorisation property of the chiral
bar complex. The chiral Chevalley--Eilenberg complex, being
a chiral coalgebra, satisfies the factorisation axiom
by construction (Beilinson--Drinfeld~\cite{BD04}, Chapter 3).
The $\Pthree$ bracket, being a $\cD$-module morphism
(Proposition~\ref{prop:chiral-p3-structure}(iii)), is
compatible with the factorisation isomorphisms.
\end{proof}


% ----------------------------------------------------------------
\subsection{The $\cD$-module structure and the KZ connection}
\label{subsec:chiral-e3-dmod}

The chiral $\Ethree$-algebra
$\CE^{\mathrm{ch}}_k(\fg)$ is not an abstract
$\Ethree$-algebra: it is an $\Ethree$-algebra object
in the category of $\cD$-modules on $X$, or equivalently
on $\Ran(X)$. This $\cD$-module structure is the content
that distinguishes the chiral construction from the
topological CFG construction.

\begin{proposition}[The $\cD$-module structure]
\label{prop:chiral-e3-dmod}
The chiral $\Ethree$-algebra
$\CE^{\mathrm{ch}}_k(\fg)$ carries a flat connection
\begin{equation}\label{eq:chiral-e3-connection}
  \nabla^{\mathrm{ch}}
  \;=\;
  d \;+\; \sum_{i < j}
  \frac{k\,\Omega_{ij}}{z_i - z_j}\,dz_{ij}
\end{equation}
on $\Conf_n(X)$, where $\Omega_{ij}$ is the Casimir element
acting on the $i$-th and $j$-th tensor factors. This is the
Knizhnik--Zamolodchikov connection.
\begin{enumerate}[label=\textup{(\roman*)}]
\item The flatness
  $(\nabla^{\mathrm{ch}})^2 = 0$ follows from the
  classical Yang--Baxter equation for the $r$-matrix
  $r(z) = k\,\Omega/z$
  \textup{(}trace-form convention; at $k = 0$, $r = 0$%
  \textup{)}: the curvature of $\nabla^{\mathrm{ch}}$
  is proportional to the LHS of the CYBE, which vanishes.
\item The $\Ethree$-algebra operations are horizontal
  with respect to $\nabla^{\mathrm{ch}}$: the cup product,
  the Gerstenhaber bracket, and the $\Pthree$ bracket are
  all flat sections of the corresponding $\Hom$-bundles.
  This is the $\cD$-linearity of
  Proposition~\textup{\ref{prop:chiral-p3-structure}(iii)}
  in the multi-point form.
\item The monodromy representation of~$\nabla^{\mathrm{ch}}$
  on $\pi_1(\Conf_n(X))$ gives the braid group action.
  For $X = \CC$, the monodromy of the KZ connection
  is computed by the Drinfeld associator
  $\Phi_{\mathrm{KZ}}$
  and recovers the braided monoidal structure on
  $\mathrm{Rep}(U_\hbar(\fg))$.
\end{enumerate}
\end{proposition}

\begin{proof}
The KZ connection on the $n$-point conformal blocks of
$V_k(\fg)$ is classical (Knizhnik--Zamolodchikov, 1984).
The presentation here identifies this connection as the
$\cD$-module structure on the chiral Chevalley--Eilenberg
complex, viewed as a factorisation algebra on $\Ran(X)$.

For (i): the curvature of $\nabla^{\mathrm{ch}}$ on
$\Conf_n(X)$ is
$[\nabla_i, \nabla_j]
= k^2\,[\Omega_{ij}/(z_i-z_j),\,
\Omega_{ik}/(z_i-z_k) + \Omega_{jk}/(z_j-z_k)]
\,dz_{ij} \wedge dz_{ik}$
plus cyclic permutations. The vanishing of this expression
is equivalent to the classical Yang--Baxter equation
$[r_{12}, r_{13}] + [r_{12}, r_{23}] + [r_{13}, r_{23}] = 0$
for $r_{ij} = k\,\Omega_{ij}/(z_i - z_j)$.

For (ii): the $\Ethree$ operations are defined in terms of
the PVA $\lambda$-bracket, which is itself a $\cD$-module
morphism. The cup product is a $\cD$-module map by the
commutative algebra structure on the factorisation algebra.
The $\Pthree$ bracket involves the derivation $\partial$
(from the $\lambda$ in the PVA bracket), which is the
$\cD$-module action; its compatibility with the connection
follows from the PVA sesquilinearity axiom.

For (iii): this is the Drinfeld--Kohno theorem
(Drinfeld~\cite{EK96}, Kohno, Kassel): the monodromy of the
KZ connection on $\Conf_n(\CC)$ factors through the pure
braid group $P_n$ and gives, via the Drinfeld associator
$\Phi_{\mathrm{KZ}} \in \exp(\hat{\mathfrak{t}}_3)$,
a braided monoidal structure on $\mathrm{Rep}(U_\hbar(\fg))$.
\end{proof}

% ----------------------------------------------------------------
\subsection{Comparison with CFG: formal disk global sections}
\label{subsec:chiral-e3-cfg-comparison}

The chiral $\Ethree$-algebra
$\CE^{\mathrm{ch}}_k(\fg)$ lives on $\Ran(X)$ and
carries the KZ connection; the CFG algebra $\cA^\lambda$
lives on $\RR^3$ and is purely topological. The formal
disk $D = \Spec \CC[[z]]$ is the bridge: restricting to
$D$ trivializes the $\cD$-module structure and recovers the
CFG algebra as the fiber.

\begin{theorem}[Formal disk restriction recovers CFG]
\label{thm:chiral-e3-cfg}
Let $D = \Spec \CC[[z]]$ be the formal disk in~$X$.
There is a natural map of filtered $\Ethree$-algebras
\begin{equation}\label{eq:cfg-comparison}
  \Gamma(D,\, \CE^{\mathrm{ch}}_k(\fg))
  \;\longrightarrow\;
  \cA^\lambda
\end{equation}
from the global sections of the chiral $\Ethree$-algebra
on the formal disk to the CFG $\Ethree$-algebra of
Theorem~\textup{\ref{thm:cfg}}, where the CFG coupling
$\lambda$ equals the departure from critical level
$k + h^\vee$.
\begin{enumerate}[label=\textup{(\roman*)}]
\item \textup{(Associated graded.)}
  At the level of associated graded algebras,
  \begin{equation}\label{eq:cfg-assoc-graded}
    \mathrm{gr}\,\Gamma(D,\, \CE^{\mathrm{ch}}_k(\fg))
    \;\simeq\;
    C^*(\fg)
    \;=\;
    \Sym(\fg^\vee[-1]),
  \end{equation}
  the Chevalley--Eilenberg cochains of $\fg$: both the
  chiral construction restricted to the formal disk and the
  CFG construction have the same classical limit.
\item \textup{(Matching of $\Pthree$ brackets.)}
  The restriction of the chiral $\Pthree$
  bracket~\eqref{eq:chiral-p3-total} to the formal disk
  recovers the CFG $\Pthree$ bracket.
  On the formal disk, the derivation $\partial = \partial_z$
  acts on $\CC[[z]]$; the global sections
  $\Gamma(D, \fg^*[-1] \otimes \omega_D)
  \cong \fg^*[-1][[z]]\,dz$ produce, after extracting the
  residue at $z = 0$, the finite-dimensional $\Pthree$
  bracket
  \[
    \{\phi_a, \phi_b\}_{\mathrm{CFG}}
    \;=\;
    f^{ab}{}_c\,\phi_c
    \;+\;
    k\,(a,b) \cdot \delta_0,
  \]
  where $\delta_0$ is the delta-function at the origin,
  the formal-disk avatar of $\partial$. Under the
  identification $\Gamma(D, \omega_D) \cong \CC\,dz$
  \textup{(}taking the constant mode\textup{)}, this
  recovers the CFG bracket on $C^*(\fg)$.
\item \textup{(Matching of deformation spaces.)}
  Both $\CE^{\mathrm{ch}}_k(\fg)\big|_D$ and
  $\cA^\lambda$ are deformations of $C^*(\fg)$ over the
  base $(k+h^\vee)\,H^3(\fg)[[k+h^\vee]]$, with the same
  one-dimensional fibre at each order in $(k+h^\vee)$ for
  simple~$\fg$.
  The map~\eqref{eq:cfg-comparison} intertwines
  the $(k+h^\vee)$-adic filtrations.
\end{enumerate}
\end{theorem}

\begin{proof}
The argument has two steps: construct the $\Ethree$-algebra
structure on $\Gamma(D, \CE^{\mathrm{ch}}_k(\fg))$
by symmetric monoidal transport, then identify it with the
CFG algebra $\cA^\lambda$.

\smallskip
\noindent\textbf{Step 1: symmetric monoidal trivialization
on~$D$.}
The formal disk $D = \Spec \CC[[z]]$ is contractible:
every $\cD$-module on~$D$ (that is, a $\CC[[z]]$-module
with flat connection) is isomorphic to its fibre at the
origin, since every flat connection on~$D$ is
trivializable ($\pi_1(D^n) = 0$ for all~$n$, so there is
no monodromy obstruction). The global sections functor
\begin{equation}\label{eq:disk-equivalence}
  \Gamma(D, -) \colon
  \cD\textrm{-}\mathrm{mod}(D) \;\xrightarrow{\;\sim\;}\;
  \mathrm{Vect}
\end{equation}
is an equivalence of categories. This equivalence is
symmetric monoidal: the tensor product of $\cD$-modules
on~$D$ (the $!$-tensor product, which on the contractible
formal disk coincides with the $*$-tensor product and with
the ordinary tensor product of the underlying vector
spaces) corresponds under~\eqref{eq:disk-equivalence} to
the tensor product in $\mathrm{Vect}$.

An equivalence of symmetric monoidal categories preserves
all algebraic structures definable in terms of the monoidal
structure. In particular, it preserves $\Ethree$-algebra
structures: if $\cB$ is an $\Ethree$-algebra in
$\cD\textrm{-}\mathrm{mod}(D)$, then
$\Gamma(D, \cB)$ is an $\Ethree$-algebra in
$\mathrm{Vect}$, with structure maps obtained by applying
$\Gamma(D, -)$ to those of~$\cB$.

Applying this to
$\cB = \CE^{\mathrm{ch}}_k(\fg)\big|_D$: the
restriction of the chiral $\Ethree$-algebra to~$D$ is an
$\Ethree$-algebra in $\cD\textrm{-}\mathrm{mod}(D)$ (by
restriction of the $\Ethree$-structure on~$X$), and the
equivalence~\eqref{eq:disk-equivalence} endows
$\Gamma(D, \CE^{\mathrm{ch}}_k(\fg))$ with an
$\Ethree$-algebra structure in $\mathrm{Vect}$. This
$\Ethree$-structure is functorial and carries the
$(k+h^\vee)$-adic filtration inherited from the chiral side.

\smallskip
\noindent\textbf{Step 2: identification with $\cA^\lambda$.}
It remains to identify
$\Gamma(D, \CE^{\mathrm{ch}}_k(\fg))$ with the CFG
algebra $\cA^\lambda$ as filtered $\Ethree$-algebras,
where $\lambda = k + h^\vee$. Both are filtered
$\Ethree$-algebras deforming the same classical limit over
the same base, and the identification proceeds through
three structural invariants.

Part (i): on~$D$, the $\cD$-module structure trivializes,
so the chiral Chevalley--Eilenberg complex reduces to the
ordinary one. The associated graded of the chiral
$\Ethree$-algebra on~$D$ is
\[
  \mathrm{gr}\,\Gamma(D,\, \CE^{\mathrm{ch}}_k(\fg))
  \;=\;
  \Gamma(D,\, \mathrm{gr}\,\CE^{\mathrm{ch}}_k(\fg))
  \;=\;
  \Gamma(D,\, \Sym^c(\fg^*[-1]) \otimes \omega_D).
\]
On the formal disk, $\omega_D \cong \cO_D\,dz$; the
equivalence~\eqref{eq:disk-equivalence} extracts the
constant mode in the $z$-expansion, giving
$\Sym(\fg^\vee[-1]) = C^*(\fg)$, the
Chevalley--Eilenberg cochains of~$\fg$ as a commutative
algebra. This is the associated graded of $\cA^\lambda$
(Theorem~\ref{thm:cfg}(i)).

Part (ii): the chiral $\Pthree$ bracket on~$X$
(equation~\eqref{eq:chiral-p3-total}) is defined via the
PVA $\lambda$-bracket. On the formal disk, the spectral
parameter $\lambda$ becomes the formal variable of the
completed enveloping algebra: the translation generator
$\partial_z$ acts on $\CC[[z]]$, and taking global sections
via~\eqref{eq:disk-equivalence} extracts the residue at
$z = 0$. The resulting bracket on
$\Gamma(D, \fg^*[-1] \otimes \omega_D)
\cong \fg^*[-1][[z]]\,dz$ is
\[
  \{\phi_a, \phi_b\}_{\mathrm{CFG}}
  \;=\;
  f^{ab}{}_{c}\,\phi_c
  \;+\;
  k\,(a,b) \cdot \delta_0,
\]
where $f^{ab}{}_{c}$ are the structure constants of~$\fg$,
$(a,b)$ is the Killing form, and $\delta_0$ is the
delta-function at the origin (the formal-disk avatar of the
derivation~$\partial$). Under the identification
$\Gamma(D, \omega_D) \cong \CC\,dz$, this is the CFG
$\Pthree$ bracket on $C^*(\fg)$
(Theorem~\ref{thm:cfg}(ii)).

Part (iii): the deformation theory of $\Ethree$-algebras
with fixed associated graded $C^*(\fg)$ is controlled by
the $\Ethree$-Hochschild cohomology, which for
$C^*(\fg)$ with~$\fg$ simple is
$(k+h^\vee)\,H^3(\fg)[[k+h^\vee]]$, one-dimensional at each
order in $(k+h^\vee)$. Both
$\Gamma(D, \CE^{\mathrm{ch}}_k(\fg))$ and
$\cA^\lambda$ are deformations of $C^*(\fg)$ over this
base (Theorem~\ref{thm:e3-cs}(ii) and
Theorem~\ref{thm:cfg}(iv) respectively). Since the
deformation space is one-dimensional at each order, and
both deformations have matching associated graded
(Part~(i)) and matching $\Pthree$ bracket (Part~(ii)),
they are isomorphic as filtered $\Ethree$-algebras.
The $(k+h^\vee)$-adic filtration on the chiral side maps to
the $\lambda$-adic filtration on the CFG side under
$\lambda = k + h^\vee$.
\end{proof}

\begin{remark}[The chain of spaces]
\label{rem:disk-vs-point}
The restriction passes through
$\mathrm{pt} \leftarrow D \to \mathbb{A}^1 \to \PP^1
\to X$, with each step adding geometric content: the formal
disk provides the derivation $\partial$ (cocycle term
$k(a,b)\partial$); $\mathbb{A}^1$ adds Arnold relations
and braid monodromy; $\PP^1$ and higher-genus $X$ add
compactness and modular curvature. None of these
restrictions is an identity; each step loses geometric
information.
\end{remark}


% ----------------------------------------------------------------
\subsection{The Khan--Zeng topological enhancement}
\label{subsec:khan-zeng-enhancement}

The chiral $\Ethree$-algebra is holomorphic: it depends on the
complex structure of $X$ through the KZ connection. At
non-critical level $k \neq -h^\vee$, the Sugawara Virasoro
element provides a homotopy between the holomorphic and
topological directions, upgrading the $\Ethree$ structure
from holomorphic to topological. At the critical level,
the Sugawara element is undefined and the enhancement fails.

\begin{proposition}[Topological enhancement via Sugawara]
\label{prop:khan-zeng-topological}
At generic non-critical level $k \neq -h^\vee$, the affine
Kac--Moody vertex algebra $V_k(\fg)$ possesses a Sugawara
Virasoro element
\begin{equation}\label{eq:sugawara-element}
  T_{\mathrm{Sug}} \;=\;
  \frac{1}{2(k + h^\vee)}
  \sum_{a=1}^{\dim \fg}
  \mathopen{:} J^a J_a \mathclose{:},
\end{equation}
which generates an inner action of the Virasoro algebra at
central charge $c = k\,\dim(\fg)/(k + h^\vee)$.

By Khan--Zeng~\cite{KhanZeng25}, the Virasoro element
upgrades the holomorphic-topological symmetry of the $3$d
gauge theory on $\Sigma \times \RR_t$ to fully topological
symmetry on $\Sigma \times \RR_t$. At the level of the
chiral $\Ethree$-algebra, this means:
\begin{enumerate}[label=\textup{(\roman*)}]
\item \textup{(Holomorphic-topological.)}
  The chiral $\Ethree$-algebra
  $\CE^{\mathrm{ch}}_k(\fg)$ lives on $X \times \RR$,
  where $X$ carries holomorphic structure and $\RR$
  carries topological structure. The $\Etwo$ from the
  curve geometry is holomorphic; the $+1$ from the $\Pthree$
  bracket is topological \textup{(}from the $\RR$ factor%
  \textup{)}.
\item \textup{(Full topological.)}
  The Sugawara element provides an identification
  of the holomorphic direction in~$X$ with a topological
  direction: the stress-energy tensor $T_{\mathrm{Sug}}$
  generates translations along $X$, making them homotopically
  trivial. After this identification, the
  $\Ethree$-structure becomes \emph{topological} $\Ethree$:
  an $\Ethree$-algebra in the category of locally constant
  sheaves on $\RR^3 \cong X \times \RR$
  \textup{(}identifying $X \cong \CC \cong \RR^2$%
  \textup{)}.
\item \textup{(Failure at critical level.)}
  At $k = -h^\vee$, the Sugawara element~\eqref{eq:sugawara-element}
  is undefined \textup{(}the denominator $2(k+h^\vee)$
  vanishes\textup{)}. The critical-level chiral
  $\Ethree$-algebra $\CE^{\mathrm{ch}}_{-h^\vee}(\fg)$ is
  genuinely holomorphic, not topological: the
  $\Etwo$-coalgebra on $\BarSig(\cA)$ is topological
  (Warning~\ref{warn:e1-vs-e2-source}), but $\cA$
  itself retains essential dependence on the complex
  structure of~$X$.
  Without the Sugawara element, there is no homotopy
  between the holomorphic and topological directions,
  so the full structure is not topological $\Ethree$.
  Perturbatively, in the formal neighbourhood of critical level,
  the topological enhancement holds order-by-order in
  $(k + h^\vee)$ \textup{(}since $T_{\mathrm{Sug}}$ exists at
  any nonzero departure from critical level\textup{)}.
\end{enumerate}
\end{proposition}

\begin{proof}
Part (i) is the HT structure of the $3$d gauge theory
whose boundary chiral algebra is $V_k(\fg)$
(Costello--Gwilliam~\cite{CG17}). The factorisation
algebra of the HT theory lives on $\Sigma \times \RR$
with holomorphic dependence on $\Sigma$ and topological
dependence on $\RR$; the chiral $\Ethree$-algebra is the
algebraic encoding of this factorisation structure.

Part (ii).
Khan--Zeng~\cite{KhanZeng25} prove the following at the
classical level: an inner Virasoro element in a Poisson
vertex algebra with non-degenerate central charge provides
a homotopy between the holomorphic and topological
directions, upgrading the HT theory to a fully topological
theory. We extend this to the quantum vertex algebra
$V_k(\fg)$. The Sugawara construction
(Frenkel--Ben-Zvi~\cite{FBZ04}, Chapter~15) produces
$T_{\mathrm{Sug}} \in V_k(\fg)$ satisfying the exact
Virasoro OPE
\[
  T_{\mathrm{Sug}}(z)\,T_{\mathrm{Sug}}(w) \;\sim\;
  \frac{c/2}{(z-w)^4}
  + \frac{2\,T_{\mathrm{Sug}}(w)}{(z-w)^2}
  + \frac{\partial T_{\mathrm{Sug}}(w)}{z-w}
\]
at central charge $c = k\,\dim(\fg)/(k+h^\vee)$, and the
Sugawara OPE with the current generators
\[
  T_{\mathrm{Sug}}(z)\,J^a(w) \;\sim\;
  \frac{J^a(w)}{(z-w)^2}
  + \frac{\partial J^a(w)}{z-w}.
\]
The second relation gives $L_{-1} \cdot J^a = \partial J^a$:
the holomorphic translation operator $L_{-1} = (T_{\mathrm{Sug}})_{(1)}$
acts as $\partial$ on all generators. These are identities in
the quantum vertex algebra, not formal-deformation statements;
the normal ordering in~\eqref{eq:sugawara-element} absorbs
all quantum corrections.  The BRST operator $Q$ of the $3$d
gauge theory satisfies $[Q, b] = L_{-1}$ where $b$ is the
antighost zero-mode (this is the standard BRST relation for
the Sugawara Virasoro; see~\cite{FBZ04}, \S15.4), so
$L_{-1}$ is BRST-exact. The Khan--Zeng
argument~\cite{KhanZeng25} then applies verbatim: BRST-exactness
of holomorphic translation is the datum required to upgrade
the HT symmetry to full topological symmetry. At generic
$k$, the central charge
$c = k\,\dim(\fg)/(k+h^\vee) \neq 0$, so the hypothesis is
satisfied.

Part (iii): at $k = -h^\vee$, the Feigin--Frenkel centre
$Z(V_{-h^\vee}(\fg)) = \mathrm{Fun}(\mathrm{Op}_{\fg^\vee}(D))$
(Feigin--Frenkel) is a commutative
algebra with non-trivial dependence on the complex structure
of the formal disk (through the notion of $\fg^\vee$-oper).
The absence of the Sugawara element means no homotopy
between holomorphic and topological directions exists, and
the chiral $\Ethree$-algebra is genuinely holomorphic.
The perturbative statement follows because the Sugawara
element exists as a formal power series in $(k + h^\vee)$ for
$k \neq -h^\vee$.
\end{proof}


% ----------------------------------------------------------------
\subsection{Example: $\mathfrak{sl}_2$ chiral $\Ethree$-algebra}
\label{subsec:sl2-chiral-e3}

The general construction specialises cleanly to
$\fg = \mathfrak{sl}_2$: the chiral $\Pthree$ bracket is
determined by six pairings on the generators
$\phi_e, \phi_f, \phi_h$, each consisting of a Lie bracket
term (from the structure constants) and a cocycle term
(proportional to the level $k$, vanishing at $k = 0$).
Restricting to the formal disk and extracting the zero mode
recovers the CFG bracket of
Proposition~\ref{prop:e3-explicit-sl2}.

\begin{example}[Explicit chiral $\Pthree$ bracket for
$\mathfrak{sl}_2$]
\label{ex:sl2-chiral-e3}
For $\fg = \mathfrak{sl}_2$ with basis $\{e, f, h\}$,
Killing form $(e,f) = (f,e) = 1$, $(h,h) = 2$,
structure constants $[e,f] = h$, $[h,e] = 2e$, $[h,f] = -2f$,
and $h^\vee = 2$, the chiral $\Pthree$ bracket
(Construction~\ref{constr:chiral-p3-bracket}) on the
generators $\phi_e, \phi_f, \phi_h \in \mathfrak{sl}_2^*[-1]$
is:
\begin{align}
  \label{eq:sl2-p3-ef}
  \{\phi_e, \phi_f\}^{\mathrm{ch}}
  &= \phi_h + k\,\partial, \\
  \label{eq:sl2-p3-he}
  \{\phi_h, \phi_e\}^{\mathrm{ch}}
  &= 2\,\phi_e, \\
  \label{eq:sl2-p3-hf}
  \{\phi_h, \phi_f\}^{\mathrm{ch}}
  &= -2\,\phi_f, \\
  \label{eq:sl2-p3-hh}
  \{\phi_h, \phi_h\}^{\mathrm{ch}}
  &= 2k\,\partial, \\
  \label{eq:sl2-p3-ee}
  \{\phi_e, \phi_e\}^{\mathrm{ch}}
  &= 0, \\
  \label{eq:sl2-p3-ff}
  \{\phi_f, \phi_f\}^{\mathrm{ch}}
  &= 0.
\end{align}
The Lie bracket terms ($\phi_h$ in~\eqref{eq:sl2-p3-ef},
$2\phi_e$ in~\eqref{eq:sl2-p3-he},
$-2\phi_f$ in~\eqref{eq:sl2-p3-hf}) persist at $k = 0$;
the cocycle terms ($k\,\partial$ in~\eqref{eq:sl2-p3-ef}
and~\eqref{eq:sl2-p3-hh}) vanish at $k = 0$
(AP126 check: at $k = 0$, the level-dependent terms
vanish, and the bracket reduces to the Chevalley--Eilenberg
coderivation).

At the critical level $k = -2$ (i.e.\ $k + h^\vee = 0$), the cocycle
terms become $-2\,\partial$, and the chiral
$\Ethree$-algebra degenerates: the curvature $m_0 = 0$
(uncurved), and the commutative factorisation structure is
undeformed.

Taking formal disk global sections and extracting the zero
mode recovers the CFG $\Pthree$ bracket of
Proposition~\ref{prop:e3-explicit-sl2}:
$\{\phi_a, \phi_b\}_{\mathrm{CFG}}
= f^{ab}{}_c\,\phi_c + k\,(a,b)\,\delta_0$, which at the
cohomological level gives
$\{X, Y\} = h_{\mathrm{KZ}}\,(X,Y)$ for
$X, Y \in \HH^1 = \mathfrak{sl}_2$
(equation~\eqref{eq:e3-p3-killing}), with
$h_{\mathrm{KZ}} = 1/(k+2)$.
\end{example}

% ================================================================
\section{Examples}
\label{sec:examples}
